\documentclass[a4paper]{article}

%%%%%%%%%%%%%%%%
%%% PREAMBLE %%%
%%%%%%%%%%%%%%%%

%%% PACKAGES %%%
\usepackage{fontspec}                     % set fonts
	\setmainfont{Junicode}
\usepackage[a4paper,margin=3cm]{geometry} % page layout
\usepackage{fancyhdr}                     % headers & footers
\usepackage{titling}                      % macros \thetitle, \theauthor
\usepackage{graphicx}                     % enhanced graphics support
	\graphicspath{{img/}}
\usepackage[automake,toc]{glossaries}     % glossaries (obv)
	\makeglossaries
	\newglossaryentry{turnstile}
{%
	name=turnstile,
	description={An ingress control device to help with the organizization
	of the passengers’ boarding. Basically a lockable revolving gate with
	three arms. When signalled, the gate unlocks and a single passenger may
	go through by turning one of the arms. After a single rotation, the gate
	locks again, waiting for the next signal.}
}

\newglossaryentry{authentication}
{%
	name=authentication,
	description={The process of verifying the \textit{identity} of the
	passengers.}
}

\newglossaryentry{authorization}
{%
	name=authorization,
	description={The process of veryfying whether an identified passenger
	has \textit{permission} to do something (eg board the bus).}
}

\newglossaryentry{alarm}
{%
	name=alarm,
	description={A device installed at every door that emits light and sound
	to draw attention to certain events. For instance, it signals the end of
	the boarding phase and also turns on in the event of an unauthorized
	boarding attempt.}
}

\newglossaryentry{RFIDScanner}
{%
	name={RFID scanner},
	description={An instrument capable of reading radio frequency
	identification tags. This can be used to effectively identify
	passengers.}
}

\newglossaryentry{boardingController}
{%
	name={boarding controller},
	description={A device responsible for the orchestration of the whole
	boarding phase (signalling doors to open or close, activating the alarm,
	etc).}
}

\usepackage{enumitem}                     % enhanced enumerations
\usepackage{tabularx}

%%% OTHER %%%
\setlength\headheight{22.3725pt}

%%% META %%%
\title{Assignment \#1 \\ Requirement analysis}
\author{Levendula}
\date{\today}

%%% ADDITIONAL PACKAGE CONFIG %%%
% fancyhdr
\pagestyle{fancy}
\fancyhf{}
\lhead{\theauthor}
\rhead{\thetitle}
\lfoot{\today}
\rfoot{\thepage}

%%%%%%%%%%%%%%%%%%%%%%%%%%%%%%%%%%%%%%%%%%%%%%%%%%%%%%%%%%%%%%%%%%%%%%%%%%%%%%%%

%%%%%%%%%%%%
%%% BODY %%%
%%%%%%%%%%%%

\begin{document}

\begin{titlepage}
	\begin{center}
		\includegraphics[width=8cm]{logo.jpg}

		\vspace{.2cm}

		\textbf{Budapest University of Technology and Economics} \\
		Faculty of Electrical Engineering and Informatics \\
		Department of Measurement and Information Systems \\

		\vspace{2cm}

		{\huge IT System Design (\texttt{VIMIAC01})}

		\vspace{2cm}

		{\huge \bfseries \thetitle}

		\vspace{.5cm}

		{\Large \theauthor}

		\vspace{.5cm}

		{\Large \today}
	\end{center}

	\vfill{}

	{\large Team members:}

	\vspace{.25cm}

	\begin{tabular}{lll}
		Annamária Gálik &
			\texttt{WGMUO2} &
			ancsi666@gmail.com \\
		Borbála Szilágyi &
			\texttt{COVQ1M} &
			szilagyiborbala8@gmail.com \\
		Bertalan Z. Péter &
			\texttt{QO7CU6} &
			bertalan.peter+uni@bertalanp99.eu
	\end{tabular}

	\vspace{2cm}
\end{titlepage}


\tableofcontents
\clearpage

\section{Our task}

% TODO basically a copy of the task description and/or the specification

% Ancsi: Nem kell a dolgokat {} közé tenni, csak ritkán (akkor, ha alkalmazni
% akarsz rájuk valami olyan makrót, ami ilyen blokkokra teljesül, pl \scriptsize
% (de ilyet csak kivételes esetekben csinálunk)
% A sorok végére sem kell tenni \\-t, hanem egy sort kihagyva érjük el, hogy új
% paragrafus legyen
% Az ilyen TODO jellegű dolgokat pedig kommentbe tegyük (%-kal kezdődő sor)
% -- Berci

Our job was to do the preparation of installing driverless autonomous vehicles
to offer transportation service to the workers at a private office park.

We had to identify the stakeholders and model the context of the planned system,
and also to define the high-level requirements, necessary functions and use
cases. We also documented non-trivial terms in a glossary as well as the
ambiguous or contradicting requirements so that we can contact the customer with
well-defined questions.

We had to study the appendices, then we had to look through and discuss the
questions that had arisen.

% TODO: The goal of the project???


\section{Ambiguous requirements}

\begin{enumerate}
	\item Interpretation of ‘make sure the buses are not overcrowded’

		Ideas:

		\textit{The buses form a network that is capable of
			intelligently deciding which bus to ‘send’ where, in
			accordance to the number of people requesting the bus at
			the terminals, what routes other buses have planned and
			how many people are on certain buses. This approach is
			to solve the problem mentioned above; overcrowding can
			be avoided provided there is a sufficient number of
			vehicles at our disposal and the buses can autonomously
			decide their routes in such a way that no single bus is
			full while others are almost empty. Essentially, we
			could balance the number of passengers on the buses.}

	\item What does a request contain of? Who can publish a request?

		Ideas:

		\textit{The request must contains the indentification of the
			worker who ordered the bus, the number of passengers who
			will travel, the starting point, the destination, and
			the planned time of the trip. To publish a request
			needed an authentication to avoid unauthorized access.}

	\item Can the passenger change his/her destination during the
		transportation?

		Ideas:

		\textit{No, but the passengers can get off the vehicle before
			their former destination if they indicate their
			intention in time.}
\end{enumerate}


\section{List of stakeholders}

% Ancsi: Erre van egy külön environment: description
%
% Úgy működik, hogy a \item-ek után írsz egy [] közé valamit, ami a magyarázandó
% szó, majd utána írod is a magyarázatot. pl:
%
% \begin{description}
% 	\item[HTTP] HyperText Transfer Protocol
%	\item[MPLS] MultiProtocol Label Switching
% 	\item[SMTP] Simple Mail Transfer Protocol
% \end{description}
%
% (itt ennél kicsit mást használtam, hogy még szebb legyen, de az egy package
% által működik csak (enumitem))
%
% Persze lehet szebbben formázni, ahogy lent tettem, mert hosszúak a sorok.
%
% -- Berci

\begin{description}[align=right,leftmargin=6cm,style=multiline]
	\item[authorities]
		Experts of autonomous vehicles with official responsibility for
		a particular area of activity
		% TODO ez szerintem nem ez, hanem a hatóságok -- Berci

	\item[office workers]
		The people working at the client's company and the main users of
		the autonomous buses.

	\item[Health\&Safety Department]
		A group at our company responsible for analyzing the
		potential health and safety risks of our products and
		forming rules and regulations to make sure these risks
		are averted.

	\item[project sponsor]
		An organization that contributes financially to the project.

	\item[project manager]
		A person who takes responsibility for the planning, preparation
		and execution of the project. They are responsible for
		accomplishing the project's objectives.

	\item[architects]
		The people who design new architectures. They are responsible
		for achieving a particular plan or aim.
		% TODO az architect ilyen építész, nem általános emberke
		% (gondolom itt csak a buszmegállóra vontakozik)

	\item[autonomous vehicle control engineers]
		They are capable of developing new technologies and handling
		problems of autonomous vehicle transport systems taking into
		account environmental and energy management requirements.

	\item[mechanical engineers]
		They work with automotive mechanical systems, and design
		power-producing machines.
		% TODO power-producing machines?

	\item[mobile application developers]
		They create, maintain and implement mobile applications that
		meet the requirements of the clients.

	\item[terminal developers]
		% TODO one sentence a One sentence about their duty. I don't
		% really know right now, sorry.

	\item[system testers]
		They create testing plans for the system and identify what parts
		of a system can be tested using automated tools and what require
		manual testing. This team is also responsible for running the
		tests.

	\item[mechanics]
		People who repair and maintain machinery.

	\item[suppliers]
		The organizations that provide needed elements, products and
		services for the project.

	\item[our company’s management]
		The group of people in charge of our company.

	\item[client's management]
		The group of people in charge of our client's company.
\end{description}



\section{System context}

% TODO


\section{Requirements}

\subsection{Functional requirements}
\begin{tabularx}{\textwidth}{p{.75cm} X}
	F1     & The buses shall get around autonomously inside the site. \\
	F1.1   & The buses shall be capable of transporting passengers around
	         the office park. \\
	F1.2   & The buses shall select their route automatically based on the
	         current demand. \\
        F1.3   & The buses must not have any permanent personnel on board. \\
	F2     & The buses shall have a dedicated station in the park for
	         storage and maintenance. \\
	F3     & The passengers should be able to publish their transportation
	         requests. \\
	F3.1   & Registered users should be able to publish their requests in
	         advance using a mobile application. \\
	F3.2   & Passengers should be able to publish their transportation
	         requests through the terminals. \\
	F3.2.1 & Dedicated terminals should be installed in front of every major
	         office building. \\
\end{tabularx}

% TODO: kep beszurasa a diagramrol

\subsection{Performance requirements}
\begin{tabularx}{\textwidth}{p{.75cm} X}
        P1 & The waiting time for a bus may not exceed 30 minutes. \\
        P2 & The buses shall operate between 7 AM and 10 PM every workday. \\
\end{tabularx}

% TODO: kep beszurasa a diagramrol

\subsection{Reliability requirements}
\begin{tabularx}{\textwidth}{p{.75cm} X}
	R1 & The buses shall be able to complete at least 100 km on average
	     between two interventions. \\
	R2 & At least ⅔ of the buses should be operational at all times. \\
        R3 & The buses should balance their load evenly if possible. \\
\end{tabularx}

% TODO: kep beszurasa a diagramrol

\subsection{Safety requirements}
\begin{tabularx}{\textwidth}{p{.75cm} X}
        SA01 & The buses must avoid causing any harm to the passengers. \\
	SA02 & The system should have a database which contains actual
	       information about the condition of the buses. \\
        SA03 & The buses must avoid causing accidents. \\
	SA04 & The buses should be equipped with sensors that can monitor road
	       conditions. \\
	SA05 & The buses must successfully complete a 30-day test run before
	       acceptance. \\
        SA06 & The buses may not be overcrowded. \\
        SA07 & Health\&Safety regulations must be fulfilled. \\
        SA08 & The buses may have a black box. \\
	SA09 & The buses should close their doors only when no passenger is
	       standing in the doorway. \\
	SA10 & The buses should only brake gracefully, except when abrupt
	       braking is the only viable option to avoid an accident. \\
	SA11 & The Health\&Safety Department should have access to the database.
	       \\
	SA12 & The autonomous vehicles must be able to react to the behaviour of
	       cyclists and pedestrians. \\
\end{tabularx}

% TODO: kep beszurasa a diagramrol

\subsection{Security requirements}
\begin{tabularx}{\textwidth}{p{.75cm} X}
	SE1 & The system must provide control access only authorized personnel.
	      \\
	SE2 & The mobile application must protect the personal information of
	      the users. \\
\end{tabularx}

% TODO: kep beszurasa a diagramrol

\subsection{Supportability requirements}
\begin{tabularx}{\textwidth}{p{.75cm} X}
	SU1 & The buses should be able to notify the maintenance team and return
	      to the station autonomously in case of a non-critical failure. \\
	SU2 & Buses may have a group of 2–3 people stationed on site who can
	      intervene (either remotely or manually) in case of a technical
	      problem. \\
\end{tabularx}

% TODO: kep beszurasa a diagramrol

\subsection{Usability requirements}
\begin{tabularx}{\textwidth}{p{.75cm} X}
        U1 & The buses must be able to operate under any weather conditions. \\
        U2 & The bus stations may be covered. \\
	U3 & The mobile application’s login screen should be simple and
	     efficient. \\
\end{tabularx}

% TODO: kep beszurasa a diagramrol


\section{Use cases}

\begin{description}[style=nextline]
	\item[Make request via mobile application]
		Use the mobile application designed especially for this system
		to make a request for a lift. Only registered users may make
		requests this way. The request includes the user's
		identification, the date and time of the travel, their starting
		and destination point as well as the number of people who wish
		to travel in case of a group request.

	\item[Make request via terminal]
		Use the installed terminals in front of every office building to
		make a tranportation request. When making a request this way,
		only the travel destionation can (and must) be specified.

	\item[Monitor system status]
		Periodically check if the entire system is operating as
		expected — ie an appropriate number of buses is in operation
		(with regard to the number of users at the moment), the buses
		calculate routes correctly and do not develop any faults. It
		must also be checked that the platforms for making requests can
		be published are fully operational.

	\item[Monitor bus condition]
		Keep watch on the various parameters of a bus to make sure it is
		functioning adequately.

	\item[Monitor bus position]
		Keep track of the positions of the buses to make sure they are
		located where they are supposed to be at any given time.

	\item[Avoid causing accidents]
		Ensure the buses cause harm to neither their environment nor
		their passengers. This is largely accomplished by collecting
		data from a variety of sensors and reacting in every situation
		intelligently and handling unexpected occurrences.

	\item[Drive bus]
		Control the movement of a bus in order to fulfill the main
		purpose of the system: to transport passengers.

	\item[Find alternative route]
		When a bus cannot run on a given route (which can occur due to a
		variety of reasons) it must find an alternative route to reach
		its destination provided such a path exists. The buses are aware
		of the road layout of the office park.

	\item[Select route]
		The act of choosing a path between two locations in the office
		park, preferably one that is the shortest possible.

	\item[Transport passengers]
		Get passengers to their desired destinations.

	\item[Manage requests]
		Take requests from users and use them to route the buses in such
		a way that the requests are carried out efficiently and quickly.

	\item[Monitor environment]
		The buses are aware of their surrounding area using several
		sensors installed for this purpose.

	\item[Notify mechanics]
		In case of a fault the bus ‘calls’ for repairs. Repairs can be
		made either remotely or manually. When possible, the buses
		return to the nearest terminal where such repairs can be made.

	\item[Repair and maintain]
		Fix any potential problems with a bus and keep it in shape for
		continued usage.

	\item[Travel]
		Use the system to get from point A to point B.
\end{description}


\section{Work journal}

% TODO copy from Sheets when done

\clearpage
\printglossaries % TODO does not work

\end{document}
